\documentclass[11pt]{article}

\usepackage[margin=0.75cm]{geometry}
\usepackage{amsmath}
\usepackage{amssymb}
\usepackage{enumitem}
\usepackage{listings}
\usepackage[T1]{fontenc}

\setlength{\columnsep}{0.7cm}
\setlength{\parindent}{0pt}
\setlist{nosep,leftmargin=*}

\lstset{
  basicstyle=\ttfamily\scriptsize,
  columns=fullflexible,
  keepspaces=true,
  aboveskip=2pt,
  belowskip=2pt,
  frame=single,
  breaklines=true
}

\begin{document}
\twocolumn

\begin{center}
{\Large \textbf{Problem Set 7: Classes and Recursion}}\\[-1mm]
COSC 120 -- Prof. Jordan C. Hanson
\end{center}

\textbf{Goal:} Review the class concepts from Chapter 11 and the
recursive patterns from Chapter 12 using problems closely related to
our warm-up and in-class exercises.

\textbf{Submit:} Your Python files and a PDF containing the requested
written answers, traces, and explanations.

\section*{1. A 2D Particle Class}

Create a module named

\begin{lstlisting}
particle2d.py
\end{lstlisting}

containing a class

\begin{lstlisting}
class Particle2D:
\end{lstlisting}

with the four instance attributes

\[
x,\qquad y,\qquad v_x,\qquad v_y.
\]

Your constructor should have the form

\begin{lstlisting}
def __init__(self, x, y, vx, vy):
    ...
\end{lstlisting}

Add the following methods.

\textbf{A. Update the particle state}

\begin{lstlisting}
def update(self, dt):
    ...
\end{lstlisting}

Use

\[
x_{\rm new}=x+v_x\Delta t,
\qquad
y_{\rm new}=y+v_y\Delta t.
\]

The method should change the object's position but leave its velocity
unchanged.

\textbf{B. Return the particle speed}

\begin{lstlisting}
def speed(self):
    ...
\end{lstlisting}

using

\[
v=\sqrt{v_x^2+v_y^2}.
\]

\textbf{C. Return the distance from the origin}

Add

\begin{lstlisting}
def distance(self):
    ...
\end{lstlisting}

using

\[
r=\sqrt{x^2+y^2}.
\]

Create a second file,

\begin{lstlisting}
test_particle.py
\end{lstlisting}

that imports the class and creates

\begin{lstlisting}
p1 = Particle2D(1, 2, 3, -1)
p2 = Particle2D(-2, 0, 1, 2)
\end{lstlisting}

For each particle:

\begin{enumerate}
\item print its initial position, speed, and distance from the origin;
\item call \texttt{update(0.5)};
\item print its new position and distance.
\end{enumerate}

Explain briefly why updating \texttt{p1} does not change \texttt{p2}.

\section*{2. Recursion with a List}

Write a recursive function

\begin{lstlisting}
def all_nonnegative(values):
    ...
\end{lstlisting}

that returns \texttt{True} only if every number in the list is greater
than or equal to zero.

Your function should follow the same recursive pattern used in class:

\begin{itemize}
\item the \textbf{base case} occurs when no values remain;
\item examine one element using \texttt{pop()};
\item if that element fails the test, stop immediately;
\item otherwise call the same function on the smaller list.
\end{itemize}

Test the function using

\begin{lstlisting}
a = [4, 9, 0, 12, 3, 7]
b = [4, 9, -2, 12, 3, 7]
\end{lstlisting}

Before running your code, predict the result for each list.

Then answer:

\begin{enumerate}
\item What is the base case?
\item How does each recursive call make the problem smaller?
\item Why can the function stop early when a negative number is found?
\end{enumerate}

\section*{3. Recursive Factorial}

Write

\begin{lstlisting}
def factorial(n):
    ...
\end{lstlisting}

using recursion.

Your function should use

\[
0! = 1
\]

as the base case and

\[
n! = n(n-1)!
\]

as the recursive case.

\textbf{A.} Compute and print

\begin{lstlisting}
factorial(7)
\end{lstlisting}

\textbf{B.} Show the recursive expansion of \(5!\), beginning with

\[
5! = 5(4!)
\]

and continuing until the base case is reached.

\textbf{C.} Add a conditional so that a negative input such as

\begin{lstlisting}
factorial(-3)
\end{lstlisting}

does not recurse forever. Your function may return \texttt{None} for a
negative input.

Explain why a recursive function needs a base case.

\section*{4. Recursive Binary Search}

Use the sorted list

\begin{lstlisting}
values = [2, 6, 9, 13, 18,
          21, 27, 31, 36, 42,
          48, 55, 63, 72, 84]
\end{lstlisting}

Write a recursive binary-search function

\begin{lstlisting}
def binary_search(values, target,
                  low, high):
    ...
\end{lstlisting}

that keeps the original list and changes only the indices
\texttt{low} and \texttt{high}.

Use

\[
\texttt{mid}
=
\left\lfloor
\frac{\texttt{low}+\texttt{high}}{2}
\right\rfloor.
\]

Your function must:

\begin{itemize}
\item return \texttt{-1} when \texttt{low > high};
\item return \texttt{mid} when the target is found;
\item recursively search the left half when the target is smaller;
\item recursively search the right half when the target is larger.
\end{itemize}

\textbf{A. Search for a value that is present.}

Call

\begin{lstlisting}
binary_search(
    values, 63, 0, len(values)-1)
\end{lstlisting}

and record the index returned.

Complete a trace with columns

\[
\texttt{low},\quad
\texttt{high},\quad
\texttt{mid},\quad
\texttt{values[mid]},\quad
\texttt{keep}.
\]

In the \texttt{keep} column, write \textbf{left}, \textbf{right}, or
\textbf{found}.

\textbf{B. Search for a missing value.}

Call

\begin{lstlisting}
binary_search(
    values, 50, 0, len(values)-1)
\end{lstlisting}

and record the sequence of midpoint values examined.

What should the function return?

\textbf{C. Explain the algorithm.}

In two or three sentences, explain why binary search requires the list
to be sorted and why it is more efficient than checking every element
one at a time for a large sorted list.

\section*{5. Files to Submit}

Submit:

\begin{lstlisting}
particle2d.py
test_particle.py
recursion_tools.py
problem_set_7.pdf
\end{lstlisting}

Place \texttt{all\_nonnegative()}, \texttt{factorial()}, and
\texttt{binary\_search()} in \texttt{recursion\_tools.py}.

Your PDF should contain the requested predictions, recursive traces,
and short explanations.

\end{document}
